% titel: Blatt 3, A9: Enthalpieänderungen von Ethanol
% Teil II: der durchgerechnete Weg.

\subsection*{Gegeben und gesucht}
$m(\mathrm{C_2H_6O}) = \qty{10}{g}$, $M \approx \qty{46.07}{g/mol}$;
$c_p = \qty{2.42}{J\,g^{-1}K^{-1}}$,
$\Delta_{vap}H^\circ = \qty{43.5}{kJ/mol}$,
$\Delta_cH^\circ = \qty{-1368}{kJ/mol}$,
$\Delta_fH^\circ(\text{flüssig}) = \qty{-277.69}{kJ/mol}$.
Gesucht: Enthalpieänderung für a) Erwärmen 20--40\,°C, b) Verdampfen,
c) Verbrennen, d) Bilden aus den Elementen.

\subsection*{Ansatz}
Gemeinsame Stoffmenge $n = m/M$; a) über $m\,c_p\,\Delta T$, b)--d) über
$n\cdot$ molare Enthalpie.

\subsection*{Rechnung}
\begin{equation*}
  n = \frac{10}{46.07} = 0.2171\ \mathrm{mol}.
\end{equation*}
\begin{itemize}
  \item[a)] $\Delta T = 40 - 20 = \qty{20}{K}$:
        \begin{equation*}
          \Delta H_a = 10\times 2.42 \times 20 = \qty{484}{J}
          = \qty{0.484}{kJ}.
        \end{equation*}
  \item[b)] $\Delta H_b = n\,\Delta_{vap}H^\circ
        = 0.2171\times 43.5 = \qty{9.44}{kJ}$.
  \item[c)] $\Delta H_c = n\,\Delta_cH^\circ
        = 0.2171\times(-1368) = \qty{-297}{kJ}$.
  \item[d)] $\Delta H_d = n\,\Delta_fH^\circ
        = 0.2171\times(-277.69) = \qty{-60.3}{kJ}$.
\end{itemize}

\begin{ergebnis}
\begin{align*}
  \text{a) Erwärmen } &\Delta H \approx \qty{0.48}{kJ},\\
  \text{b) Verdampfen } &\Delta H \approx \qty{9.4}{kJ},\\
  \text{c) Verbrennen } &\Delta H \approx \qty{-297}{kJ},\\
  \text{d) Bilden } &\Delta H \approx \qty{-60.3}{kJ}.
\end{align*}
\end{ergebnis}

\begin{kontrolle}
Wärmemengen/Grössenordnungen plausibel: Erwärmen kleiner Betrag
(0.5 kJ), Verdampfen positiv und grösser, Verbrennen dominiert (stark
exotherm, ca. 30 kJ/g entsprechend der üblichen Verbrennungswärme von
Ethanol). Querprobe d) $\to$ c): $\Delta_cH$ liegt wegen der
Verbrennungsprodukte weiter negativ als $\Delta_fH$; beides mit demselben
$n = 0.217\ \mathrm{mol}$ gerechnet.
\end{kontrolle}

\begin{merksatz}
Erwärmen mit $m\,c_p\Delta T$, Phasen- und Reaktionsenthalpien mit
$n\cdot\Delta H^\circ$ rechnen; Vorzeichen von Verdampfung (positiv) und
Verbrennung (negativ) beherzigen.
\end{merksatz}
